\documentclass[11pt]{article}
\usepackage[a4paper,margin=1in]{geometry}
\usepackage{amsmath,amssymb,mathtools}
\usepackage{booktabs}
\usepackage{enumitem}
\usepackage[bookmarks=false,hidelinks]{hyperref}
\usepackage{microtype}
\allowdisplaybreaks

\newcommand{\F}{\mathbb{F}}
\newcommand{\Z}{\mathbb{Z}}
\newcommand{\SpanA}{\langle a\rangle}
\newcommand{\mymod}{\pmod m}


\title{Supplementary Material for\\
\emph{On the Minimum Distance of High-rate Moore-Type Quasi-Cyclic LDPC Codes}\\[0.5ex]
\large Detailed Algebraic Derivations for Weight-$8$ Cases II--VII}
\author{}
\date{}

\begin{document}
\maketitle

\section{Scope and relation to the main manuscript}
This supplementary document provides the detailed algebraic derivations omitted from the main manuscript for Cases II--VII in the characterization of Hamming-weight-$8$ codewords of $\mathrm{Moore}(m,a,b,3)$ under the assumption $\operatorname{ord}_m(b)=3$.  Case I is derived in full in the main manuscript; the purpose of the present document is to supply the corresponding calculations for the remaining six matching configurations.  The notation, normalization, case labels, and condition labels \emph{(a)}--\emph{(d)} are those of the main manuscript.

The correspondence between the case labels used below and the matching subcases is
\begin{center}
\begin{tabular}{@{}cccc@{}}
\toprule
Case & Matching subcase & RREF & Free variables \\
\midrule
II   & $1\text{-}2\text{-}1$ & $A_{121}$ & $u_4,u_7,u_8$ \\
III  & $1\text{-}2\text{-}2$ & $A_{122}$ & $u_6,u_7,u_8$ \\
IV   & $2\text{-}1\text{-}2$ & $A_{212}$ & $u_4,u_6,u_7,u_8$ \\
V    & $2\text{-}1\text{-}7$ & $A_{217}$ & $u_4,u_6,u_7,u_8$ \\
VI   & $2\text{-}3\text{-}5$ & $A_{235}$ & $u_4,u_6,u_7,u_8$ \\
VII  & $2\text{-}4\text{-}3$ & $A_{243}$ & $u_5,u_6,u_7,u_8$ \\
\bottomrule
\end{tabular}
\end{center}

\section{Algebraic framework and conventions}
All calculations are carried out in the finite field $\F_m$.  We use the normalization $c_1=(0,0)$ and write
\[
u_i=a^{t_i},\qquad u_1=1.
\]
Since $\operatorname{ord}_m(b)=3$, we have $b\neq1$ and
\[
1+b+b^2=0.
\]
All coefficients below are simplified using this identity.

For a matching edge $\{c_i,c_q\}$ at level $s\in\{0,1,2\}$, the basic relation is
\[
x_i-x_q=(b^s-1)(u_q-u_i).
\]
For an oriented cycle
\[
\underline{i_1,i_2,\ldots,i_\ell}(s_1s_2\cdots s_\ell),
\]
the underlined indices denote the successive vertices and $s_r$ is the level of the edge joining $c_{i_r}$ to $c_{i_{r+1}}$, with $i_{\ell+1}=i_1$.  Summing the basic edge relations around the cycle eliminates the $x$-coordinates and gives a homogeneous linear equation in $u_1,\ldots,u_8$.

For each matching configuration we proceed as follows.  First, we choose a spanning tree and form its fundamental-cycle system.  Second, we row-reduce the resulting homogeneous system in the fixed column order $(u_1,\ldots,u_8)$; zero rows are omitted.  Third, after imposing $u_1=1$, the first pivot relation gives condition \emph{(b)}, while the remaining pivot-variable expressions give condition \emph{(d)} through membership in $\SpanA$.  Finally, the level-$0$ coordinates and within-block distinctness conditions are reduced, modulo the main congruence, to conditions \emph{(a)} and \emph{(c)}.  Factors that differ by a nonzero scalar in $\F_m$ are identified, since they define the same nonvanishing condition.

\paragraph{Labeling convention for Case VI.}
For subcase $2\text{-}3\text{-}5$, we use the representative appearing in the proposition of the main manuscript.  Relative to the intermediate labeling arising during the enumeration, this representative is obtained by the relabeling $c_7\leftrightarrow c_8$.  This convention fixes the column order and yields the matrix $A_{235}$ reported in the main manuscript.

\section{Case II: Subcase $1\text{-}2\text{-}1$}
\subsection{Matching configuration and fundamental-cycle system}
The matching set is
\[
\begin{aligned}
M_0&=\left\{\{c_{1},c_{2}\},\ \{c_{3},c_{4}\},\ \{c_{5},c_{6}\},\ \{c_{7},c_{8}\}\right\},\\
M_1&=\left\{\{c_{1},c_{3},c_{6},c_{8}\},\ \{c_{2},c_{5}\},\ \{c_{4},c_{7}\}\right\},\\
M_2&=\left\{\{c_{1},c_{4}\},\ \{c_{2},c_{3}\},\ \{c_{5},c_{8}\},\ \{c_{6},c_{7}\}\right\}.
\end{aligned}
\]
We choose the spanning tree
\[
F={\{c_{1},c_{3}\}_{1},\ \{c_{1},c_{6}\}_{1},\ \{c_{1},c_{8}\}_{1},\ \{c_{2},c_{5}\}_{1},\ \{c_{4},c_{7}\}_{1},\ \{c_{1},c_{4}\}_{2},\ \{c_{2},c_{3}\}_{2}}.
\]
This tree has $7$ edges.  Since the matching graph has $16$ edges, it determines $9$ fundamental cycles.  The associated cycle equations are listed below.
\begin{enumerate}[label=(\arabic*)]
\item $ \underline{3,1,6}(111) $.  The resulting cycle equation is
\[
\begin{aligned}
&0=0.
\end{aligned}
\]
\item $ \underline{3,1,8}(111) $.  The resulting cycle equation is
\[
\begin{aligned}
&0=0.
\end{aligned}
\]
\item $ \underline{6,1,8}(111) $.  The resulting cycle equation is
\[
\begin{aligned}
&0=0.
\end{aligned}
\]
\item $ \underline{5,2,3,1,8}(12112) $.  The resulting cycle equation is
\[
\begin{aligned}
&\left(1 + 2b\right)u_{2} + \left(-1 - 2b\right)u_{3} + \left(-1 - 2b\right)u_{5} + \left(1 + 2b\right)u_{8}=0.
\end{aligned}
\]
\item $ \underline{6,1,4,7}(1212) $.  The resulting cycle equation is
\[
\begin{aligned}
&\left(1 + 2b\right)u_{1} + \left(-1 - 2b\right)u_{4} + \left(-1 - 2b\right)u_{6} + \left(1 + 2b\right)u_{7}=0.
\end{aligned}
\]
\item $ \underline{1,3,2}(120) $.  The resulting cycle equation is
\[
\begin{aligned}
&\left(1 - b\right)u_{1} + \left(-2 - b\right)u_{2} + \left(1 + 2b\right)u_{3}=0.
\end{aligned}
\]
\item $ \underline{3,1,4}(120) $.  The resulting cycle equation is
\[
\begin{aligned}
&\left(1 + 2b\right)u_{1} + \left(1 - b\right)u_{3} + \left(-2 - b\right)u_{4}=0.
\end{aligned}
\]
\item $ \underline{5,2,3,1,6}(12110) $.  The resulting cycle equation is
\[
\begin{aligned}
&\left(1 + 2b\right)u_{2} + \left(-1 - 2b\right)u_{3} + \left(1 - b\right)u_{5} + \left(-1 + b\right)u_{6}=0.
\end{aligned}
\]
\item $ \underline{7,4,1,8}(1210) $.  The resulting cycle equation is
\[
\begin{aligned}
&\left(-1 - 2b\right)u_{1} + \left(1 + 2b\right)u_{4} + \left(1 - b\right)u_{7} + \left(-1 + b\right)u_{8}=0.
\end{aligned}
\]
\end{enumerate}
After deleting zero equations and row-reducing the remaining system in the column order $(u_1,\ldots,u_8)$, we obtain
\[
A_{121}=
\begin{bmatrix}
1 & 0 & 0 & -1 & 0 & 0 & 1 + b & -1 - b \\
0 & 1 & 0 & -1 & 0 & 0 & 2 + b & -2 - b \\
0 & 0 & 1 & -1 & 0 & 0 & 1 & -1 \\
0 & 0 & 0 & 0 & 1 & 0 & 1 + b & -2 - b \\
0 & 0 & 0 & 0 & 0 & 1 & b & -1 - b
\end{bmatrix}.
\]
Thus the fundamental-cycle system has the RREF reported in the main manuscript.

\subsection{RREF solution and conditions \emph{(b)} and \emph{(d)}}
The free variables are
\[
u_{4},\ u_{7},\ u_{8}.
\]
The pivot variables are therefore
\[
\begin{aligned}
u_{1}&=u_{4} + \left(-1 - b\right)u_{7} + \left(1 + b\right)u_{8},\\
u_{2}&=u_{4} + \left(-2 - b\right)u_{7} + \left(2 + b\right)u_{8},\\
u_{3}&=u_{4} - u_{7} + u_{8},\\
u_{5}&=\left(-1 - b\right)u_{7} + \left(2 + b\right)u_{8},\\
u_{6}&=-bu_{7} + \left(1 + b\right)u_{8}.
\end{aligned}
\]
Imposing $u_1=1$ gives the main congruence
\[
-1 + u_{4} + \left(-1 - b\right)u_{7} + \left(1 + b\right)u_{8}=0,
\]
which is condition \emph{(b)} after substituting $u_i=a^{t_i}$.  Requiring the remaining dependent variables above to belong to $\SpanA$ gives condition \emph{(d)}.

\subsection{Distinctness conditions \emph{(a)} and \emph{(c)}}
For the distinctness calculation we solve the main congruence for $u_{4}$:
\[
u_{4}=1 + \left(1 + b\right)u_{7} + \left(-1 - b\right)u_{8}.
\]
After this substitution, the common $x$-coordinate of each block of $M_0$ is
\[
\begin{aligned}
x_{1}=x_{2}&=0,\\
x_{3}=x_{4}&=\left(1 + 2b\right)u_{7} + \left(-1 - 2b\right)u_{8},\\
x_{5}=x_{6}&=-1 + b + \left(-1 - 2b\right)u_{7} + \left(2 + b\right)u_{8},\\
x_{7}=x_{8}&=-1 + b + \left(1 - b\right)u_{8}.
\end{aligned}
\]
Hence different level-$0$ blocks have distinct $x$-coordinates, and the columns inside each block are distinct, exactly when the following factors are all nonzero (factors differing by a nonzero scalar are identified):
\[
\begin{gathered}
u_{7} - u_{8}\not\equiv 0\mymod,\\
1 + bu_{7} + \left(-1 - b\right)u_{8}\not\equiv 0\mymod,\\
1 - u_{8}\not\equiv 0\mymod,\\
1 + 2bu_{7} + \left(-1 - 2b\right)u_{8}\not\equiv 0\mymod.
\end{gathered}
\]
Using the main congruence, these factors reduce to condition \emph{(a)} together with the nonzero factors in condition \emph{(c)}.  Explicitly,
\[
\begin{aligned}
\text{a1: }\ t_7 \neq t_8&\Longleftrightarrow u_{7} - u_{8}\not\equiv0\mymod,\\
\text{a2: }\ t_8 \neq 0&\Longleftrightarrow 1 - u_{8}\not\equiv0\mymod,\\
\text{c1}&\Longleftrightarrow 1 + bu_{7} + \left(-1 - b\right)u_{8}\not\equiv0\mymod,\\
\text{c2}&\Longleftrightarrow 1 + 2bu_{7} + \left(-1 - 2b\right)u_{8}\not\equiv0\mymod.
\end{aligned}
\]
Hence conditions \emph{(a)} and \emph{(c)} are precisely the level-$0$ and within-block distinctness requirements for this matching configuration.  Together with conditions \emph{(b)} and \emph{(d)}, they give the necessary and sufficient algebraic conditions for this case.
\section{Case III: Subcase $1\text{-}2\text{-}2$}
\subsection{Matching configuration and fundamental-cycle system}
The matching set is
\[
\begin{aligned}
M_0&=\left\{\{c_{1},c_{2}\},\ \{c_{3},c_{4},c_{5},c_{6}\},\ \{c_{7},c_{8}\}\right\},\\
M_1&=\left\{\{c_{1},c_{3}\},\ \{c_{2},c_{5}\},\ \{c_{4},c_{7}\},\ \{c_{6},c_{8}\}\right\},\\
M_2&=\left\{\{c_{1},c_{4}\},\ \{c_{2},c_{6}\},\ \{c_{3},c_{7}\},\ \{c_{5},c_{8}\}\right\}.
\end{aligned}
\]
We choose the spanning tree
\[
F={\{c_{1},c_{3}\}_{1},\ \{c_{2},c_{5}\}_{1},\ \{c_{4},c_{7}\}_{1},\ \{c_{6},c_{8}\}_{1},\ \{c_{1},c_{4}\}_{2},\ \{c_{2},c_{6}\}_{2},\ \{c_{1},c_{2}\}_{0}}.
\]
This tree has $7$ edges.  Since the matching graph has $16$ edges, it determines $9$ fundamental cycles.  The associated cycle equations are listed below.
\begin{enumerate}[label=(\arabic*)]
\item $ \underline{3,1,4,7}(1212) $.  The resulting cycle equation is
\[
\begin{aligned}
&\left(1 + 2b\right)u_{1} + \left(-1 - 2b\right)u_{3} + \left(-1 - 2b\right)u_{4} + \left(1 + 2b\right)u_{7}=0.
\end{aligned}
\]
\item $ \underline{5,2,6,8}(1212) $.  The resulting cycle equation is
\[
\begin{aligned}
&\left(1 + 2b\right)u_{2} + \left(-1 - 2b\right)u_{5} + \left(-1 - 2b\right)u_{6} + \left(1 + 2b\right)u_{8}=0.
\end{aligned}
\]
\item $ \underline{3,1,4}(120) $.  The resulting cycle equation is
\[
\begin{aligned}
&\left(1 + 2b\right)u_{1} + \left(1 - b\right)u_{3} + \left(-2 - b\right)u_{4}=0.
\end{aligned}
\]
\item $ \underline{3,1,2,5}(1010) $.  The resulting cycle equation is
\[
\begin{aligned}
&\left(-1 + b\right)u_{1} + \left(1 - b\right)u_{2} + \left(1 - b\right)u_{3} + \left(-1 + b\right)u_{5}=0.
\end{aligned}
\]
\item $ \underline{3,1,2,6}(1020) $.  The resulting cycle equation is
\[
\begin{aligned}
&\left(-1 + b\right)u_{1} + \left(2 + b\right)u_{2} + \left(1 - b\right)u_{3} + \left(-2 - b\right)u_{6}=0.
\end{aligned}
\]
\item $ \underline{4,1,2,5}(2010) $.  The resulting cycle equation is
\[
\begin{aligned}
&\left(-2 - b\right)u_{1} + \left(1 - b\right)u_{2} + \left(2 + b\right)u_{4} + \left(-1 + b\right)u_{5}=0.
\end{aligned}
\]
\item $ \underline{4,1,2,6}(2020) $.  The resulting cycle equation is
\[
\begin{aligned}
&\left(-2 - b\right)u_{1} + \left(2 + b\right)u_{2} + \left(2 + b\right)u_{4} + \left(-2 - b\right)u_{6}=0.
\end{aligned}
\]
\item $ \underline{5,2,6}(120) $.  The resulting cycle equation is
\[
\begin{aligned}
&\left(1 + 2b\right)u_{2} + \left(1 - b\right)u_{5} + \left(-2 - b\right)u_{6}=0.
\end{aligned}
\]
\item $ \underline{7,4,1,2,6,8}(120210) $.  The resulting cycle equation is
\[
\begin{aligned}
&\left(-2 - b\right)u_{1} + \left(2 + b\right)u_{2} + \left(1 + 2b\right)u_{4} + \left(-1 - 2b\right)u_{6}\\ &\qquad  + \left(1 - b\right)u_{7} + \left(-1 + b\right)u_{8}=0.
\end{aligned}
\]
\end{enumerate}
After deleting zero equations and row-reducing the remaining system in the column order $(u_1,\ldots,u_8)$, we obtain
\[
A_{122}=
\begin{bmatrix}
1 & 0 & 0 & 0 & 0 & -1 + b & -1 & 1 - b \\
0 & 1 & 0 & 0 & 0 & -1 + b & 0 & -b \\
0 & 0 & 1 & 0 & 0 & b & -1 & -b \\
0 & 0 & 0 & 1 & 0 & -1 & -1 & 1 \\
0 & 0 & 0 & 0 & 1 & b & 0 & -1 - b
\end{bmatrix}.
\]
Thus the fundamental-cycle system has the RREF reported in the main manuscript.

\subsection{RREF solution and conditions \emph{(b)} and \emph{(d)}}
The free variables are
\[
u_{6},\ u_{7},\ u_{8}.
\]
The pivot variables are therefore
\[
\begin{aligned}
u_{1}&=\left(1 - b\right)u_{6} + u_{7} + \left(-1 + b\right)u_{8},\\
u_{2}&=\left(1 - b\right)u_{6} + bu_{8},\\
u_{3}&=-bu_{6} + u_{7} + bu_{8},\\
u_{4}&=u_{6} + u_{7} - u_{8},\\
u_{5}&=-bu_{6} + \left(1 + b\right)u_{8}.
\end{aligned}
\]
Imposing $u_1=1$ gives the main congruence
\[
-1 + \left(1 - b\right)u_{6} + u_{7} + \left(-1 + b\right)u_{8}=0,
\]
which is condition \emph{(b)} after substituting $u_i=a^{t_i}$.  Requiring the remaining dependent variables above to belong to $\SpanA$ gives condition \emph{(d)}.

\subsection{Distinctness conditions \emph{(a)} and \emph{(c)}}
For the distinctness calculation we solve the main congruence for $u_{7}$:
\[
u_{7}=1 + \left(-1 + b\right)u_{6} + \left(1 - b\right)u_{8}.
\]
After this substitution, the common $x$-coordinate of each block of $M_0$ is
\[
\begin{aligned}
x_{1}=x_{2}&=0,\\
x_{3}=x_{4}=x_{5}=x_{6}&=\left(-1 + b\right)u_{6} + \left(1 - b\right)u_{8},\\
x_{7}=x_{8}&=\left(-2 + 2b\right)u_{6} + \left(2 - 2b\right)u_{8}.
\end{aligned}
\]
Hence different level-$0$ blocks have distinct $x$-coordinates, and the columns inside each block are distinct, exactly when the following factors are all nonzero (factors differing by a nonzero scalar are identified):
\[
\begin{gathered}
u_{6} - u_{8}\not\equiv 0\mymod,\\
1 + \left(-1 + b\right)u_{6} - bu_{8}\not\equiv 0\mymod,\\
1 - 2u_{6} + u_{8}\not\equiv 0\mymod,\\
1 + 2bu_{6} + \left(-1 - 2b\right)u_{8}\not\equiv 0\mymod.
\end{gathered}
\]
Using the main congruence, these factors reduce to condition \emph{(a)} together with the nonzero factors in condition \emph{(c)}.  Explicitly,
\[
\begin{aligned}
\text{a1: }\ t_6 \neq t_8&\Longleftrightarrow u_{6} - u_{8}\not\equiv0\mymod,\\
\text{a2: }\ t_7 \neq t_8&\Longleftrightarrow 1 + \left(-1 + b\right)u_{6} - bu_{8}\not\equiv0\mymod,\\
\text{c1}&\Longleftrightarrow 1 - 2u_{6} + u_{8}\not\equiv0\mymod,\\
\text{c2}&\Longleftrightarrow 1 + 2bu_{6} + \left(-1 - 2b\right)u_{8}\not\equiv0\mymod.
\end{aligned}
\]
Hence conditions \emph{(a)} and \emph{(c)} are precisely the level-$0$ and within-block distinctness requirements for this matching configuration.  Together with conditions \emph{(b)} and \emph{(d)}, they give the necessary and sufficient algebraic conditions for this case.
\section{Case IV: Subcase $2\text{-}1\text{-}2$}
\subsection{Matching configuration and fundamental-cycle system}
The matching set is
\[
\begin{aligned}
M_0&=\left\{\{c_{1},c_{2}\},\ \{c_{3},c_{4}\},\ \{c_{5},c_{6}\},\ \{c_{7},c_{8}\}\right\},\\
M_1&=\left\{\{c_{1},c_{3}\},\ \{c_{2},c_{4}\},\ \{c_{5},c_{7}\},\ \{c_{6},c_{8}\}\right\},\\
M_2&=\left\{\{c_{1},c_{5}\},\ \{c_{3},c_{6}\},\ \{c_{2},c_{7}\},\ \{c_{4},c_{8}\}\right\}.
\end{aligned}
\]
We choose the spanning tree
\[
F={\{c_{1},c_{3}\}_{1},\ \{c_{2},c_{4}\}_{1},\ \{c_{5},c_{7}\}_{1},\ \{c_{6},c_{8}\}_{1},\ \{c_{1},c_{5}\}_{2},\ \{c_{2},c_{7}\}_{2},\ \{c_{3},c_{6}\}_{2}}.
\]
This tree has $7$ edges.  Since the matching graph has $12$ edges, it determines $5$ fundamental cycles.  The associated cycle equations are listed below.
\begin{enumerate}[label=(\arabic*)]
\item $ \underline{4,2,7,5,1,3,6,8}(12121212) $.  The resulting cycle equation is
\[
\begin{aligned}
&\left(-1 - 2b\right)u_{1} + \left(1 + 2b\right)u_{2} + \left(1 + 2b\right)u_{3} + \left(-1 - 2b\right)u_{4}\\ &\qquad  + \left(1 + 2b\right)u_{5} + \left(-1 - 2b\right)u_{6} + \left(-1 - 2b\right)u_{7} + \left(1 + 2b\right)u_{8}=0.
\end{aligned}
\]
\item $ \underline{1,5,7,2}(2120) $.  The resulting cycle equation is
\[
\begin{aligned}
&\left(2 + b\right)u_{1} + \left(-2 - b\right)u_{2} + \left(-1 - 2b\right)u_{5} + \left(1 + 2b\right)u_{7}=0.
\end{aligned}
\]
\item $ \underline{3,1,5,7,2,4}(121210) $.  The resulting cycle equation is
\[
\begin{aligned}
&\left(1 + 2b\right)u_{1} + \left(-1 - 2b\right)u_{2} + \left(1 - b\right)u_{3} + \left(-1 + b\right)u_{4}\\ &\qquad  + \left(-1 - 2b\right)u_{5} + \left(1 + 2b\right)u_{7}=0.
\end{aligned}
\]
\item $ \underline{5,1,3,6}(2120) $.  The resulting cycle equation is
\[
\begin{aligned}
&\left(-1 - 2b\right)u_{1} + \left(1 + 2b\right)u_{3} + \left(2 + b\right)u_{5} + \left(-2 - b\right)u_{6}=0.
\end{aligned}
\]
\item $ \underline{7,5,1,3,6,8}(121210) $.  The resulting cycle equation is
\[
\begin{aligned}
&\left(-1 - 2b\right)u_{1} + \left(1 + 2b\right)u_{3} + \left(1 + 2b\right)u_{5} + \left(-1 - 2b\right)u_{6}\\ &\qquad  + \left(1 - b\right)u_{7} + \left(-1 + b\right)u_{8}=0.
\end{aligned}
\]
\end{enumerate}
After deleting zero equations and row-reducing the remaining system in the column order $(u_1,\ldots,u_8)$, we obtain
\[
A_{212}=
\begin{bmatrix}
1 & 0 & 0 & -1 & 0 & -1 - b & b & 1 \\
0 & 1 & 0 & -1 & 0 & 0 & b & -b \\
0 & 0 & 1 & -1 & 0 & -1 - b & 0 & 1 + b \\
0 & 0 & 0 & 0 & 1 & -1 & -1 & 1
\end{bmatrix}.
\]
Thus the fundamental-cycle system has the RREF reported in the main manuscript.

\subsection{RREF solution and conditions \emph{(b)} and \emph{(d)}}
The free variables are
\[
u_{4},\ u_{6},\ u_{7},\ u_{8}.
\]
The pivot variables are therefore
\[
\begin{aligned}
u_{1}&=u_{4} + \left(1 + b\right)u_{6} - bu_{7} - u_{8},\\
u_{2}&=u_{4} - bu_{7} + bu_{8},\\
u_{3}&=u_{4} + \left(1 + b\right)u_{6} + \left(-1 - b\right)u_{8},\\
u_{5}&=u_{6} + u_{7} - u_{8}.
\end{aligned}
\]
Imposing $u_1=1$ gives the main congruence
\[
-1 + u_{4} + \left(1 + b\right)u_{6} - bu_{7} - u_{8}=0,
\]
which is condition \emph{(b)} after substituting $u_i=a^{t_i}$.  Requiring the remaining dependent variables above to belong to $\SpanA$ gives condition \emph{(d)}.

\subsection{Distinctness conditions \emph{(a)} and \emph{(c)}}
For the distinctness calculation we solve the main congruence for $u_{4}$:
\[
u_{4}=1 + \left(-1 - b\right)u_{6} + bu_{7} + u_{8}.
\]
After this substitution, the common $x$-coordinate of each block of $M_0$ is
\[
\begin{aligned}
x_{1}=x_{2}&=0,\\
x_{3}=x_{4}&=\left(1 + 2b\right)u_{7} + \left(-1 - 2b\right)u_{8},\\
x_{5}=x_{6}&=-2 - b + \left(2 + b\right)u_{6} + \left(2 + b\right)u_{7} + \left(-2 - b\right)u_{8},\\
x_{7}=x_{8}&=-2 - b + \left(1 + 2b\right)u_{6} + \left(2 + b\right)u_{7} + \left(-1 - 2b\right)u_{8}.
\end{aligned}
\]
Hence different level-$0$ blocks have distinct $x$-coordinates, and the columns inside each block are distinct, exactly when the following factors are all nonzero (factors differing by a nonzero scalar are identified):
\[
\begin{gathered}
u_{7} - u_{8}\not\equiv 0\mymod,\\
1 - u_{6} - u_{7} + u_{8}\not\equiv 0\mymod,\\
1 + \left(-1 - b\right)u_{6} - u_{7} + \left(1 + b\right)u_{8}\not\equiv 0\mymod,\\
1 - u_{6} + bu_{7} - bu_{8}\not\equiv 0\mymod,\\
1 + \left(-1 - b\right)u_{6} + bu_{7}\not\equiv 0\mymod,\\
u_{6} - u_{8}\not\equiv 0\mymod.
\end{gathered}
\]
Using the main congruence, these factors reduce to condition \emph{(a)} together with the nonzero factors in condition \emph{(c)}.  Explicitly,
\[
\begin{aligned}
\text{a1: }\ t_6 \neq t_8&\Longleftrightarrow u_{6} - u_{8}\not\equiv0\mymod,\\
\text{a2: }\ t_7 \neq t_8&\Longleftrightarrow u_{7} - u_{8}\not\equiv0\mymod,\\
\text{c1}&\Longleftrightarrow 1 - u_{6} - u_{7} + u_{8}\not\equiv0\mymod,\\
\text{c2}&\Longleftrightarrow 1 + \left(-1 - b\right)u_{6} - u_{7} + \left(1 + b\right)u_{8}\not\equiv0\mymod,\\
\text{c3}&\Longleftrightarrow 1 - u_{6} + bu_{7} - bu_{8}\not\equiv0\mymod,\\
\text{c4}&\Longleftrightarrow 1 + \left(-1 - b\right)u_{6} + bu_{7}\not\equiv0\mymod.
\end{aligned}
\]
Hence conditions \emph{(a)} and \emph{(c)} are precisely the level-$0$ and within-block distinctness requirements for this matching configuration.  Together with conditions \emph{(b)} and \emph{(d)}, they give the necessary and sufficient algebraic conditions for this case.
\section{Case V: Subcase $2\text{-}1\text{-}7$}
\subsection{Matching configuration and fundamental-cycle system}
The matching set is
\[
\begin{aligned}
M_0&=\left\{\{c_{1},c_{2}\},\ \{c_{3},c_{4}\},\ \{c_{5},c_{6}\},\ \{c_{7},c_{8}\}\right\},\\
M_1&=\left\{\{c_{1},c_{3}\},\ \{c_{2},c_{4}\},\ \{c_{5},c_{7}\},\ \{c_{6},c_{8}\}\right\},\\
M_2&=\left\{\{c_{1},c_{5}\},\ \{c_{2},c_{6}\},\ \{c_{3},c_{7}\},\ \{c_{4},c_{8}\}\right\}.
\end{aligned}
\]
We choose the spanning tree
\[
F={\{c_{1},c_{3}\}_{1},\ \{c_{2},c_{4}\}_{1},\ \{c_{5},c_{7}\}_{1},\ \{c_{6},c_{8}\}_{1},\ \{c_{1},c_{5}\}_{2},\ \{c_{2},c_{6}\}_{2},\ \{c_{1},c_{2}\}_{0}}.
\]
This tree has $7$ edges.  Since the matching graph has $12$ edges, it determines $5$ fundamental cycles.  The associated cycle equations are listed below.
\begin{enumerate}[label=(\arabic*)]
\item $ \underline{3,1,5,7}(1212) $.  The resulting cycle equation is
\[
\begin{aligned}
&\left(1 + 2b\right)u_{1} + \left(-1 - 2b\right)u_{3} + \left(-1 - 2b\right)u_{5} + \left(1 + 2b\right)u_{7}=0.
\end{aligned}
\]
\item $ \underline{4,2,6,8}(1212) $.  The resulting cycle equation is
\[
\begin{aligned}
&\left(1 + 2b\right)u_{2} + \left(-1 - 2b\right)u_{4} + \left(-1 - 2b\right)u_{6} + \left(1 + 2b\right)u_{8}=0.
\end{aligned}
\]
\item $ \underline{3,1,2,4}(1010) $.  The resulting cycle equation is
\[
\begin{aligned}
&\left(-1 + b\right)u_{1} + \left(1 - b\right)u_{2} + \left(1 - b\right)u_{3} + \left(-1 + b\right)u_{4}=0.
\end{aligned}
\]
\item $ \underline{5,1,2,6}(2020) $.  The resulting cycle equation is
\[
\begin{aligned}
&\left(-2 - b\right)u_{1} + \left(2 + b\right)u_{2} + \left(2 + b\right)u_{5} + \left(-2 - b\right)u_{6}=0.
\end{aligned}
\]
\item $ \underline{7,5,1,2,6,8}(120210) $.  The resulting cycle equation is
\[
\begin{aligned}
&\left(-2 - b\right)u_{1} + \left(2 + b\right)u_{2} + \left(1 + 2b\right)u_{5} + \left(-1 - 2b\right)u_{6}\\ &\qquad  + \left(1 - b\right)u_{7} + \left(-1 + b\right)u_{8}=0.
\end{aligned}
\]
\end{enumerate}
After deleting zero equations and row-reducing the remaining system in the column order $(u_1,\ldots,u_8)$, we obtain
\[
A_{217}=
\begin{bmatrix}
1 & 0 & 0 & -1 & 0 & -1 & -1 & 2 \\
0 & 1 & 0 & -1 & 0 & -1 & 0 & 1 \\
0 & 0 & 1 & -1 & 0 & 0 & -1 & 1 \\
0 & 0 & 0 & 0 & 1 & -1 & -1 & 1
\end{bmatrix}.
\]
Thus the fundamental-cycle system has the RREF reported in the main manuscript.

\subsection{RREF solution and conditions \emph{(b)} and \emph{(d)}}
The free variables are
\[
u_{4},\ u_{6},\ u_{7},\ u_{8}.
\]
The pivot variables are therefore
\[
\begin{aligned}
u_{1}&=u_{4} + u_{6} + u_{7} - 2u_{8},\\
u_{2}&=u_{4} + u_{6} - u_{8},\\
u_{3}&=u_{4} + u_{7} - u_{8},\\
u_{5}&=u_{6} + u_{7} - u_{8}.
\end{aligned}
\]
Imposing $u_1=1$ gives the main congruence
\[
-1 + u_{4} + u_{6} + u_{7} - 2u_{8}=0,
\]
which is condition \emph{(b)} after substituting $u_i=a^{t_i}$.  Requiring the remaining dependent variables above to belong to $\SpanA$ gives condition \emph{(d)}.

\subsection{Distinctness conditions \emph{(a)} and \emph{(c)}}
For the distinctness calculation we solve the main congruence for $u_{7}$:
\[
u_{7}=1 - u_{4} - u_{6} + 2u_{8}.
\]
After this substitution, the common $x$-coordinate of each block of $M_0$ is
\[
\begin{aligned}
x_{1}=x_{2}&=0,\\
x_{3}=x_{4}&=\left(-1 + b\right)u_{6} + \left(1 - b\right)u_{8},\\
x_{5}=x_{6}&=\left(-2 - b\right)u_{4} + \left(2 + b\right)u_{8},\\
x_{7}=x_{8}&=\left(-2 - b\right)u_{4} + \left(-1 + b\right)u_{6} + 3u_{8}.
\end{aligned}
\]
Hence different level-$0$ blocks have distinct $x$-coordinates, and the columns inside each block are distinct, exactly when the following factors are all nonzero (factors differing by a nonzero scalar are identified):
\[
\begin{gathered}
u_{6} - u_{8}\not\equiv 0\mymod,\\
u_{4} - u_{8}\not\equiv 0\mymod,\\
u_{4} - bu_{6} + \left(-1 + b\right)u_{8}\not\equiv 0\mymod,\\
u_{4} + bu_{6} + \left(-1 - b\right)u_{8}\not\equiv 0\mymod,\\
1 - u_{4} - u_{6} + u_{8}\not\equiv 0\mymod.
\end{gathered}
\]
Using the main congruence, these factors reduce to condition \emph{(a)} together with the nonzero factors in condition \emph{(c)}.  Explicitly,
\[
\begin{aligned}
\text{a1: }\ t_4 \neq t_8&\Longleftrightarrow u_{4} - u_{8}\not\equiv0\mymod,\\
\text{a2: }\ t_6 \neq t_8&\Longleftrightarrow u_{6} - u_{8}\not\equiv0\mymod,\\
\text{a3: }\ t_7 \neq t_8&\Longleftrightarrow 1 - u_{4} - u_{6} + u_{8}\not\equiv0\mymod,\\
\text{c1}&\Longleftrightarrow u_{4} + bu_{6} + \left(-1 - b\right)u_{8}\not\equiv0\mymod,\\
\text{c2}&\Longleftrightarrow u_{4} - bu_{6} + \left(-1 + b\right)u_{8}\not\equiv0\mymod.
\end{aligned}
\]
Hence conditions \emph{(a)} and \emph{(c)} are precisely the level-$0$ and within-block distinctness requirements for this matching configuration.  Together with conditions \emph{(b)} and \emph{(d)}, they give the necessary and sufficient algebraic conditions for this case.
\section{Case VI: Subcase $2\text{-}3\text{-}5$}
\subsection{Matching configuration and fundamental-cycle system}
The matching set is
\[
\begin{aligned}
M_0&=\left\{\{c_{1},c_{2}\},\ \{c_{3},c_{4}\},\ \{c_{5},c_{6}\},\ \{c_{7},c_{8}\}\right\},\\
M_1&=\left\{\{c_{1},c_{3}\},\ \{c_{2},c_{6}\},\ \{c_{5},c_{7}\},\ \{c_{4},c_{8}\}\right\},\\
M_2&=\left\{\{c_{1},c_{5}\},\ \{c_{2},c_{4}\},\ \{c_{3},c_{7}\},\ \{c_{6},c_{8}\}\right\}.
\end{aligned}
\]
We choose the spanning tree
\[
F={\{c_{1},c_{3}\}_{1},\ \{c_{2},c_{6}\}_{1},\ \{c_{4},c_{8}\}_{1},\ \{c_{5},c_{7}\}_{1},\ \{c_{1},c_{5}\}_{2},\ \{c_{2},c_{4}\}_{2},\ \{c_{1},c_{2}\}_{0}}.
\]
This tree has $7$ edges.  Since the matching graph has $12$ edges, it determines $5$ fundamental cycles.  The associated cycle equations are listed below.
\begin{enumerate}[label=(\arabic*)]
\item $ \underline{3,1,5,7}(1212) $.  The resulting cycle equation is
\[
\begin{aligned}
&\left(1 + 2b\right)u_{1} + \left(-1 - 2b\right)u_{3} + \left(-1 - 2b\right)u_{5} + \left(1 + 2b\right)u_{7}=0.
\end{aligned}
\]
\item $ \underline{6,2,4,8}(1212) $.  The resulting cycle equation is
\[
\begin{aligned}
&\left(1 + 2b\right)u_{2} + \left(-1 - 2b\right)u_{4} + \left(-1 - 2b\right)u_{6} + \left(1 + 2b\right)u_{8}=0.
\end{aligned}
\]
\item $ \underline{3,1,2,4}(1020) $.  The resulting cycle equation is
\[
\begin{aligned}
&\left(-1 + b\right)u_{1} + \left(2 + b\right)u_{2} + \left(1 - b\right)u_{3} + \left(-2 - b\right)u_{4}=0.
\end{aligned}
\]
\item $ \underline{5,1,2,6}(2010) $.  The resulting cycle equation is
\[
\begin{aligned}
&\left(-2 - b\right)u_{1} + \left(1 - b\right)u_{2} + \left(2 + b\right)u_{5} + \left(-1 + b\right)u_{6}=0.
\end{aligned}
\]
\item $ \underline{7,5,1,2,4,8}(120210) $.  The resulting cycle equation is
\[
\begin{aligned}
&\left(-2 - b\right)u_{1} + \left(2 + b\right)u_{2} + \left(-1 - 2b\right)u_{4} + \left(1 + 2b\right)u_{5}\\ &\qquad  + \left(1 - b\right)u_{7} + \left(-1 + b\right)u_{8}=0.
\end{aligned}
\]
\end{enumerate}
After deleting zero equations and row-reducing the remaining system in the column order $(u_1,\ldots,u_8)$, we obtain
\[
A_{235}=
\begin{bmatrix}
1 & 0 & 0 & b & 0 & -1 - b & -1 & 1 \\
0 & 1 & 0 & -1 & 0 & -1 & 0 & 1 \\
0 & 0 & 1 & b & 0 & 0 & -1 & -b \\
0 & 0 & 0 & 0 & 1 & -1 - b & -1 & 1 + b
\end{bmatrix}.
\]
Thus the fundamental-cycle system has the RREF reported in the main manuscript.

\subsection{RREF solution and conditions \emph{(b)} and \emph{(d)}}
The free variables are
\[
u_{4},\ u_{6},\ u_{7},\ u_{8}.
\]
The pivot variables are therefore
\[
\begin{aligned}
u_{1}&=-bu_{4} + \left(1 + b\right)u_{6} + u_{7} - u_{8},\\
u_{2}&=u_{4} + u_{6} - u_{8},\\
u_{3}&=-bu_{4} + u_{7} + bu_{8},\\
u_{5}&=\left(1 + b\right)u_{6} + u_{7} + \left(-1 - b\right)u_{8}.
\end{aligned}
\]
Imposing $u_1=1$ gives the main congruence
\[
-1 - bu_{4} + \left(1 + b\right)u_{6} + u_{7} - u_{8}=0,
\]
which is condition \emph{(b)} after substituting $u_i=a^{t_i}$.  Requiring the remaining dependent variables above to belong to $\SpanA$ gives condition \emph{(d)}.

\subsection{Distinctness conditions \emph{(a)} and \emph{(c)}}
For the distinctness calculation we solve the main congruence for $u_{4}$:
\[
u_{4}=1 + b - bu_{6} + \left(-1 - b\right)u_{7} + \left(1 + b\right)u_{8}.
\]
After this substitution, the common $x$-coordinate of each block of $M_0$ is
\[
\begin{aligned}
x_{1}=x_{2}&=0,\\
x_{3}=x_{4}&=\left(-2 - b\right)u_{6} + \left(2 + b\right)u_{8},\\
x_{5}=x_{6}&=-2 - b + \left(1 + 2b\right)u_{6} + \left(2 + b\right)u_{7} + \left(-1 - 2b\right)u_{8},\\
x_{7}=x_{8}&=-2 - b + \left(-1 + b\right)u_{6} + \left(2 + b\right)u_{7} + \left(1 - b\right)u_{8}.
\end{aligned}
\]
Hence different level-$0$ blocks have distinct $x$-coordinates, and the columns inside each block are distinct, exactly when the following factors are all nonzero (factors differing by a nonzero scalar are identified):
\[
\begin{gathered}
u_{6} - u_{8}\not\equiv 0\mymod,\\
1 + \left(-1 - b\right)u_{6} - u_{7} + \left(1 + b\right)u_{8}\not\equiv 0\mymod,\\
1 - bu_{6} - u_{7} + bu_{8}\not\equiv 0\mymod,\\
1 + \left(-2 - b\right)u_{6} - u_{7} + \left(2 + b\right)u_{8}\not\equiv 0\mymod,\\
1 + \left(-2 - b\right)u_{6} + bu_{7} + u_{8}\not\equiv 0\mymod,\\
1 + \left(-1 - b\right)u_{6} + bu_{7}\not\equiv 0\mymod,\\
u_{6} + \left(-1 - b\right)u_{7} + bu_{8}\not\equiv 0\mymod,\\
u_{7} - u_{8}\not\equiv 0\mymod.
\end{gathered}
\]
Using the main congruence, these factors reduce to condition \emph{(a)} together with the nonzero factors in condition \emph{(c)}.  Explicitly,
\[
\begin{aligned}
\text{a1: }\ t_4 \neq t_8&\Longleftrightarrow 1 + \left(-1 - b\right)u_{6} - u_{7} + \left(1 + b\right)u_{8}\not\equiv0\mymod,\\
\text{a2: }\ t_6 \neq t_8&\Longleftrightarrow u_{6} - u_{8}\not\equiv0\mymod,\\
\text{a3: }\ t_7 \neq t_8&\Longleftrightarrow u_{7} - u_{8}\not\equiv0\mymod,\\
\text{c1}&\Longleftrightarrow 1 + \left(-2 - b\right)u_{6} + bu_{7} + u_{8}\not\equiv0\mymod,\\
\text{c2}&\Longleftrightarrow 1 + \left(-1 - b\right)u_{6} + bu_{7}\not\equiv0\mymod,\\
\text{c3}&\Longleftrightarrow u_{6} + \left(-1 - b\right)u_{7} + bu_{8}\not\equiv0\mymod,\\
\text{c4}&\Longleftrightarrow 1 + \left(-2 - b\right)u_{6} - u_{7} + \left(2 + b\right)u_{8}\not\equiv0\mymod,\\
\text{c5}&\Longleftrightarrow 1 - bu_{6} - u_{7} + bu_{8}\not\equiv0\mymod.
\end{aligned}
\]
Hence conditions \emph{(a)} and \emph{(c)} are precisely the level-$0$ and within-block distinctness requirements for this matching configuration.  Together with conditions \emph{(b)} and \emph{(d)}, they give the necessary and sufficient algebraic conditions for this case.
\section{Case VII: Subcase $2\text{-}4\text{-}3$}
\subsection{Matching configuration and fundamental-cycle system}
The matching set is
\[
\begin{aligned}
M_0&=\left\{\{c_{1},c_{2}\},\ \{c_{3},c_{4}\},\ \{c_{5},c_{6}\},\ \{c_{7},c_{8}\}\right\},\\
M_1&=\left\{\{c_{1},c_{3}\},\ \{c_{2},c_{7}\},\ \{c_{4},c_{5}\},\ \{c_{6},c_{8}\}\right\},\\
M_2&=\left\{\{c_{1},c_{5}\},\ \{c_{2},c_{6}\},\ \{c_{3},c_{7}\},\ \{c_{4},c_{8}\}\right\}.
\end{aligned}
\]
We choose the spanning tree
\[
F={\{c_{1},c_{3}\}_{1},\ \{c_{2},c_{7}\}_{1},\ \{c_{4},c_{5}\}_{1},\ \{c_{6},c_{8}\}_{1},\ \{c_{1},c_{5}\}_{2},\ \{c_{2},c_{6}\}_{2},\ \{c_{3},c_{7}\}_{2}}.
\]
This tree has $7$ edges.  Since the matching graph has $12$ edges, it determines $5$ fundamental cycles.  The associated cycle equations are listed below.
\begin{enumerate}[label=(\arabic*)]
\item $ \underline{4,5,1,3,7,2,6,8}(12121212) $.  The resulting cycle equation is
\[
\begin{aligned}
&\left(-1 - 2b\right)u_{1} + \left(1 + 2b\right)u_{2} + \left(1 + 2b\right)u_{3} + \left(-1 - 2b\right)u_{4}\\ &\qquad  + \left(1 + 2b\right)u_{5} + \left(-1 - 2b\right)u_{6} + \left(-1 - 2b\right)u_{7} + \left(1 + 2b\right)u_{8}=0.
\end{aligned}
\]
\item $ \underline{1,3,7,2}(1210) $.  The resulting cycle equation is
\[
\begin{aligned}
&\left(1 - b\right)u_{1} + \left(-1 + b\right)u_{2} + \left(1 + 2b\right)u_{3} + \left(-1 - 2b\right)u_{7}=0.
\end{aligned}
\]
\item $ \underline{3,1,5,4}(1210) $.  The resulting cycle equation is
\[
\begin{aligned}
&\left(1 + 2b\right)u_{1} + \left(1 - b\right)u_{3} + \left(-1 + b\right)u_{4} + \left(-1 - 2b\right)u_{5}=0.
\end{aligned}
\]
\item $ \underline{5,1,3,7,2,6}(212120) $.  The resulting cycle equation is
\[
\begin{aligned}
&\left(-1 - 2b\right)u_{1} + \left(1 + 2b\right)u_{2} + \left(1 + 2b\right)u_{3} + \left(2 + b\right)u_{5}\\ &\qquad  + \left(-2 - b\right)u_{6} + \left(-1 - 2b\right)u_{7}=0.
\end{aligned}
\]
\item $ \underline{7,2,6,8}(1210) $.  The resulting cycle equation is
\[
\begin{aligned}
&\left(1 + 2b\right)u_{2} + \left(-1 - 2b\right)u_{6} + \left(1 - b\right)u_{7} + \left(-1 + b\right)u_{8}=0.
\end{aligned}
\]
\end{enumerate}
After deleting zero equations and row-reducing the remaining system in the column order $(u_1,\ldots,u_8)$, we obtain
\[
A_{243}=
\begin{bmatrix}
1 & 0 & 0 & 0 & -1 & 0 & -1 - b & 1 + b \\
0 & 1 & 0 & 0 & 0 & -1 & -1 - b & 1 + b \\
0 & 0 & 1 & 0 & -1 - b & 1 + b & -1 & 0 \\
0 & 0 & 0 & 1 & -1 - b & 1 + b & 0 & -1
\end{bmatrix}.
\]
Thus the fundamental-cycle system has the RREF reported in the main manuscript.

\subsection{RREF solution and conditions \emph{(b)} and \emph{(d)}}
The free variables are
\[
u_{5},\ u_{6},\ u_{7},\ u_{8}.
\]
The pivot variables are therefore
\[
\begin{aligned}
u_{1}&=u_{5} + \left(1 + b\right)u_{7} + \left(-1 - b\right)u_{8},\\
u_{2}&=u_{6} + \left(1 + b\right)u_{7} + \left(-1 - b\right)u_{8},\\
u_{3}&=\left(1 + b\right)u_{5} + \left(-1 - b\right)u_{6} + u_{7},\\
u_{4}&=\left(1 + b\right)u_{5} + \left(-1 - b\right)u_{6} + u_{8}.
\end{aligned}
\]
Imposing $u_1=1$ gives the main congruence
\[
-1 + u_{5} + \left(1 + b\right)u_{7} + \left(-1 - b\right)u_{8}=0,
\]
which is condition \emph{(b)} after substituting $u_i=a^{t_i}$.  Requiring the remaining dependent variables above to belong to $\SpanA$ gives condition \emph{(d)}.

\subsection{Distinctness conditions \emph{(a)} and \emph{(c)}}
For the distinctness calculation we solve the main congruence for $u_{5}$:
\[
u_{5}=1 + \left(-1 - b\right)u_{7} + \left(1 + b\right)u_{8}.
\]
After this substitution, the common $x$-coordinate of each block of $M_0$ is
\[
\begin{aligned}
x_{1}=x_{2}&=0,\\
x_{3}=x_{4}&=1 + 2b + \left(-2 - b\right)u_{6} - 3bu_{7} + \left(1 + 2b\right)u_{8},\\
x_{5}=x_{6}&=\left(-1 - 2b\right)u_{7} + \left(1 + 2b\right)u_{8},\\
x_{7}=x_{8}&=\left(-1 + b\right)u_{6} + \left(-1 - 2b\right)u_{7} + \left(2 + b\right)u_{8}.
\end{aligned}
\]
Hence different level-$0$ blocks have distinct $x$-coordinates, and the columns inside each block are distinct, exactly when the following factors are all nonzero (factors differing by a nonzero scalar are identified):
\[
\begin{gathered}
1 + bu_{6} + \left(-2 - b\right)u_{7} + u_{8}\not\equiv 0\mymod,\\
u_{7} - u_{8}\not\equiv 0\mymod,\\
u_{6} + bu_{7} + \left(-1 - b\right)u_{8}\not\equiv 0\mymod,\\
1 + bu_{6} + \left(-1 - b\right)u_{7}\not\equiv 0\mymod,\\
1 - u_{6} + \left(-1 - b\right)u_{7} + \left(1 + b\right)u_{8}\not\equiv 0\mymod,\\
u_{6} - u_{8}\not\equiv 0\mymod.
\end{gathered}
\]
Using the main congruence, these factors reduce to condition \emph{(a)} together with the nonzero factors in condition \emph{(c)}.  Explicitly,
\[
\begin{aligned}
\text{a1: }\ t_5 \neq t_6&\Longleftrightarrow 1 - u_{6} + \left(-1 - b\right)u_{7} + \left(1 + b\right)u_{8}\not\equiv0\mymod,\\
\text{a2: }\ t_6 \neq t_8&\Longleftrightarrow u_{6} - u_{8}\not\equiv0\mymod,\\
\text{a3: }\ t_7 \neq t_8&\Longleftrightarrow u_{7} - u_{8}\not\equiv0\mymod,\\
\text{c1}&\Longleftrightarrow 1 + bu_{6} + \left(-2 - b\right)u_{7} + u_{8}\not\equiv0\mymod,\\
\text{c2}&\Longleftrightarrow u_{6} + bu_{7} + \left(-1 - b\right)u_{8}\not\equiv0\mymod,\\
\text{c3}&\Longleftrightarrow 1 + bu_{6} + \left(-1 - b\right)u_{7}\not\equiv0\mymod,\\
\text{c4}&\Longleftrightarrow 1 - u_{6} + \left(-1 - b\right)u_{7} + \left(1 + b\right)u_{8}\not\equiv0\mymod.
\end{aligned}
\]
Hence conditions \emph{(a)} and \emph{(c)} are precisely the level-$0$ and within-block distinctness requirements for this matching configuration.  Together with conditions \emph{(b)} and \emph{(d)}, they give the necessary and sufficient algebraic conditions for this case.
\section{Conclusion}
For each of Cases II--VII, the fundamental-cycle equations derived from the corresponding matching configuration reduce to the RREF stated in the main manuscript.  Under the normalization $u_1=1$, the pivot relations yield condition \emph{(b)} and the subgroup-membership requirements in condition \emph{(d)}.  The remaining level-$0$ and within-block distinctness constraints yield exactly conditions \emph{(a)} and \emph{(c)}.  Consequently, the calculations in this supplement complete the algebraic verification of the six alternatives omitted from the detailed proof of the Hamming-weight-$8$ existence theorem.

\end{document}
